% -----------------------------------------------------------
\section{Seraph}
% -----------------------------------------------------------
	\label{SERAPH}
	
% % ----------------------
% \subsection{See also}
% % ----------------------

% ----------------------
\subsection{1828 American Dictionary of the English Language} % *1828
% ----------------------
	\label{SERAPH-1828}
	
	There is an entry in the online edition, but it's empty---it has no definition.

	% \begin{compactitem}
		% \item[]
	% \end{compactitem}

% ----------------------
\subsection{Oxford English Dictionary} % *OED
% ----------------------
	\label{SERAPH-OED}
	
	% -----
	\subsubsection{Seraph}
	% -----

		\begin{compactitem}
			\item[1.a] One of the seraphim.
			\item[1.b] \textit{figurative.} A seraphic person, an `angel'.
			\item[2] \textit{Geology.} A fossil shell.
		\end{compactitem}
		
	% -----
	\subsubsection{Seraphim}
	% -----
	
		OED's etymology:
		
			\begin{quote}
				< late Latin \textit{seraphim} (Vulgate), in manuscripts often \textit{seraphin} (= Greek \gr{σεραφιμ}, \gr{σεραφειμ}, Septuagint), < Hebrew \textit{ser\={a}ph\={i}m} (only in Isaiah vi), plural of \textit{*s\={a}r\={a}ph}, which is not recorded in the Bible, unless it be identical with the formally coincident word denoting a kind of venomous serpent, which occurs as quasi-adjective or in apposition to \textit{n\={a}\d{h}\={a}sh} serpent in Numbers xxi. and Deuteronomy viii. (English Bible `fiery serpents', after Vulgate \textit{ignitos serpentes}, \textit{serpens adurens}; Septuagint \gr{οφεις θανατουντας}, \gr{οφις δακνων}), and in Isaiah xiv. 29 and xxx. 6 with the epithet `flying' (English Bible `fiery flying serpent').
		
				Some scholars assume the identity of the word occurring in Isaiah vi. with that found in the other passages. On this view the `seraphim' seen by Isaiah flying above the throne of God represent a mythic or symbolic conception which must originally have had the form of a `fiery flying serpent', though in the vision this appears considerably modified. The word \textit{s\={a}r\={a}ph}, as the name of a kind of serpent, may belong to the root \textit{s\={a}raph} to burn, in allusion to the effect of the bite (compare Greek \gr{πρηστηρ}). This etymology has given rise to a conjecture that the celestial `seraphim' originally symbolized the lightning. Of those who reject the identity of \textit{s\={a}r\={a}ph} `seraph' with \textit{s\={a}r\={a}ph} `fiery serpent', some refer the former to the root of the Arabic \textit{sharafa} to be lofty or illustrious. Phonologically this is unobjectionable, but on other grounds it is now generally abandoned. Various suggestions of non-Hebrew (Egyptian, Assyrian, etc.) etymology have been made, but have not found wide acceptance.
				 
				The Latin form \textit{seraphin}, which is found in many manuscripts of the Vulgate, and is the source of all the forms used in English down to the 16th cent. (as well as of those in the Romanic languages), coincides with the Aramaic \textit{ser\={a}ph\={i}n}, but it is very doubtful whether it is more than a scribal error or a euphonic alteration. Compare French \textit{séraphin} (\textit{serafin}, 12th cent.), Provençal \textit{serafi}, Spanish \textit{serafin}, Portuguese \textit{seraphim}, Italian \textit{serafino} (all masculine singular).
				 
				In the Latin liturgical passages from which the word first became widely known, it was probably originally apprehended correctly as a plural, and readers of the Latin Bible would be guided aright by the syntax of Isaiah vi. 2; but there is evidence that `Cherubim and Seraphim' were often supposed to be the names of two individual angels. From the 15th to the 18th cent. the English plural ending was often appended, but \textit{seraphin} as a singular = `one of the seraphim' does not appear in English till late in the 16th cent. (the form \textit{seraphim} in this use not till the 17th cent.). After the introduction (perhaps by Milton) of the form \textsc{seraph} n.1, the misuse of the plural forms in singular sense gradually became rare, and it is now obsolete.
			\end{quote}
	
		\begin{compactitem}
			\item[1] In Biblical use: The living creatures with six wings, hands and feet, and a (presumably) human voice, seen in Isaiah's vision as hovering above the throne of God.
			\item[2] By Christian interpreters the seraphim were from an early period supposed to be a class of angels, and the name, associated with that of the cherubim, was introduced in the Eucharistic preface and subsequently in the \textit{Te Deum}, and thus became extensively known. The presumed derivation of the word from a Hebrew root meaning `to burn' (see above) led to the view that the seraphim are specially distinguished by fervour of love (while the cherubim excel in knowledge), and to the symbolic use of red as the colour appropriate to the seraphim in artistic representations. In the system of the Pseudo-Dionysius, the chief source of later angelology, the seraphim are the highest, and the cherubim the second, of the nine orders of angels.
			\item[2.a] seraphin (\textit{obsolete} exc. \textit{poetic} as \textit{nonce-use}), \textit{seraphim}, used as \textit{plural}. (Some of the early examples are ambiguous, and may belong to b.)
			\item[2.b] Taken as the name of an angel.
			\item[2.c] Plural.
			\item[2.d] \textit{seraphin}, \textit{seraphim} as a singular = one of the seraphim, a seraph.
			\item[3] (Meanings in heraldry.)
			\item[4] A Swedish order of knighthood.
			\item[5] (A geology use.)
			\item[6] (A moth.)
		\end{compactitem}
	
% ----------------------
\subsection{Buck's Theological Dictionary (1833)} % *BTD
% ----------------------
	\label{SERAPH-BTD}
	
	Does not appear.

	% \begin{compactdesc}
		% \item[PAGE\#]
	% \end{compactdesc}
	
% % ----------------------
% \subsection{Restoration Lexicon} % *RL *RESTORATION LEXICON
% % ----------------------
	% \label{SERAPH-OED}

	% \begin{compactitem}
		% \item[]
	% \end{compactitem}

% % ----------------------
% \subsection{Scriptural Usage} % *SCRIPTURES
% % ----------------------
	% \label{SERAPH-SCRIPTURES}

	% % -----
	% \subsubsection{Old Testament} % *OT
	% % -----
		% \label{SERAPH-OT}
	
		% % --
		% \paragraph{KJV}
		% % --
			% \label{SERAPH-OT-KJV}
	
			% \begin{tabular}{ll}
				% Total occurrences in the OT & 
			% \end{tabular}
			
			% \begin{compactdesc}
				% \item[]
			% \end{compactdesc}
			
		% % --
		% \paragraph{Hebrew Bible}
		% % --
			% \label{SERAPH-HEBREW}
		
			% %
			% \subparagraph{\texorpdfstring{\he{sample word}}{}} % paste actual hebrew text here
			% %
				% \label{PAGA} % whatever is in step bible without periods
			
				% \begin{compactdesc}
					% \item[HALOT , PDF ] \textbf{} % Use the 1994 PDF
						% \begin{compactitem}
							% \item[1]
						% \end{compactitem}
					% \item[BDB , PDF ] \textbf{}
						% \begin{compactitem}
							% \item[1]
						% \end{compactitem}
					% \item[DCH PDF ] \textbf{} % don't include the actual page number, they repeat from volume to volume
						% \begin{compactitem}
							% \item[1]
						% \end{compactitem}
				% \end{compactdesc}
				
				% \begin{compactdesc}
					% \item[Some OT uses] \textbf{}
						% \begin{compactdesc}
							% \item[]
						% \end{compactdesc}
				% \end{compactdesc}
			
		% % --
		% \paragraph{Septuagint}
		% % --
			% \label{SERAPH-LXX}
		
			% %
			% \subparagraph{\texorpdfstring{\gr{sample word}}{}}
			% %
		
	% % -----
	% \subsubsection{New Testament} % *NT
	% % -----
		% \label{SERAPH-NT}
	
		% % --
		% \paragraph{KJV}
		% % --
			% \label{SERAPH-NT-KJV}		
	
			% \begin{tabular}{ll}
				% Total occurrences in the NT & 
			% \end{tabular}
			
			% \begin{compactdesc}
				% \item[]
			% \end{compactdesc}
			
		% % --
		% \paragraph{Greek New Testament} % *GREEK
		% % --
			% \label{SERAPH-GREEK}
		
			% %
			% \subparagraph{\texorpdfstring{\gr{sample word}}{}}
			% %
				% \label{KATALLAGE}
			
				% \begin{compactdesc}
					% \item[BDAG ] \textbf{}
						% \begin{compactitem}
							% \item 
						% \end{compactitem}
					% \item[CGL , PDF ] \textbf{}
						% \begin{compactitem}
							% \item 
						% \end{compactitem}
					% \item[LSJ , PDF ] \textbf{}
						% \begin{compactitem}
							% \item 
						% \end{compactitem}
					% \item[Brill , PDF ] \textbf{}
						% \begin{compactitem}
							% \item 
						% \end{compactitem}
				% \end{compactdesc}
				
				% \begin{tabular}{ll}
					% Total occurrences in the NT & 
				% \end{tabular}
				
				% \begin{compactdesc}
					% \item[Some NT uses] \textbf{}
						% \begin{compactdesc}
							% \item[]
						% \end{compactdesc}
				% \end{compactdesc}
				
			% %
			% \subparagraph{TDNT: \texorpdfstring{\gr{sample word}}{}  (3:536--556)}
			% %
				% \label{KATALLAGE-TDNT}
		
	% % -----
	% \subsubsection{Book of Mormon} % *BOM
	% % -----
		% \label{SERAPH-BOM}
	
		% \begin{tabular}{ll}
			% Total occurrences in the BOM & 
		% \end{tabular}
		
		% \begin{compactdesc}
			% \item[]
		% \end{compactdesc}
		
	% % -----
	% \subsubsection{Doctrine and Covenants} % *DC
	% % -----
		% \label{SERAPH-DC}
	
		% \begin{tabular}{ll}
			% Total occurrences in the DC & 
		% \end{tabular}
		
		% \begin{compactdesc}
			% \item[]
		% \end{compactdesc}
		
	% % -----
	% \subsubsection{Pearl of Great Price} % *PGP
	% % -----
		% \label{SERAPH-PGP}
	
		% \begin{tabular}{ll}
			% Total occurrences in the PGP & 
		% \end{tabular}
		
		% \begin{compactdesc}
			% \item[]
		% \end{compactdesc}
		
% % ----------------------
% \subsection{Key Phrases} % *KEYPHRASES *PHRASES
% % ----------------------
	% \label{SERAPH-PHRASES}

	% \begin{compactdesc}
		% \item[] \textbf{\#times} | list
			% % \begin{compactdesc}
				% % \item[Bible anchor verses] \phantom{.}
					% % \begin{compactdesc}
						% % \item[Romans 5:5] \gr{}
					% % \end{compactdesc}
				% % \item[Commentary] ---
			% % \end{compactdesc}
	% \end{compactdesc}
	
% % ----------------------
% \subsection{Bible Dictionaries} % *DICTIONARIES
% % ----------------------
	% \label{SERAPH-BD}

% % ----------------------
% \subsection{Restoration Usage} % *RESTORATION
% % ----------------------
	% \label{SERAPH-RESTORATION}

	% % -----
	% \subsubsection{Joseph Smith} % *JS
	% % -----
		% \label{SERAPH-JS}
	
		% \begin{compactdesc}
			% \item[{\hyperref[KEY]{Date/Source}}]
		% \end{compactdesc}
		
	% % -----
	% \subsubsection{Temple Ordinances}
	% % -----
		% \label{SERAPH-TEMPLE}
	
		% \begin{compactdesc}
			% \item[{\hyperref[KEY]{Date/Source}}]
		% \end{compactdesc}
	
	% % -----
	% \subsubsection{Brigham Young} % *BY
	% % -----
		% \label{SERAPH-BY}
	
		% \begin{compactdesc}
			% \item[{\hyperref[KEY]{Date/Source}}]
		% \end{compactdesc}
	
% % ----------------------
% \subsection{Commentary and Studies} % *COMMENTARY *STUDIES
% % ----------------------
	% \label{SERAPH-COMMENTARY}

	% % -----
	% \subsubsection{Key Points}
	% % -----
		% \label{SERAPH-KEYS}
	
		% \begin{compactitem}
			% \item
		% \end{compactitem}
		
	% % -----
	% \subsubsection{Sample Study}
	% % -----
		% \label{SERAPH-study}